\documentclass[UTF8,fontset=windows,11pt,a4paper]{ctexart}
\usepackage{amsmath,amssymb,amsthm,booktabs,geometry,hyperref}
\DeclareFontShape{OMS}{cmsy}{m}{n}{<5-9>cmsy8<9->cmsy10}{}
\geometry{left=27mm,right=27mm,top=25mm,bottom=25mm}
\hypersetup{colorlinks=true,linkcolor=black,citecolor=black,urlcolor=blue}
\setlength{\parskip}{3pt}
\setlength{\emergencystretch}{2em}
\linespread{1.18}
\clubpenalty=10000\widowpenalty=10000
\pagestyle{plain}
\ctexset{section={format=\large\bfseries},subsection={format=\normalsize\bfseries}}
\newtheorem{theorem}{定理}
\renewcommand{\proofname}{证明}
\numberwithin{equation}{section}
\makeatletter
\renewcommand{\@maketitle}{\newpage\null\vskip0em
\begin{center}{\LARGE\heiti\@title\par}\vskip1em
{\normalsize\@date\par}\end{center}\par\vskip0.7em}
\makeatother
\begin{document}
\title{付费核验与经营资金：信息反馈何时减少启动资本}
\author{}\date{2026年11月1日 · 报价参数推广修订}
\maketitle
\begin{abstract}
投资者评估一项外部信号时，核验费用会减少随后能够投入经营的资金。本文研究在固定检验可靠性、期望利润和逐笔付款要求下，查看核验结果的时机以及剩余资金的用途怎样影响最低启动资本。我们先明确外部核验服务的假设，再建立参数免费已知的经营对照，将检查费用和错误退出连接到同一市场的利润损失。保证快、慢类样本的价格系数允许为任意固定正数。在信息价格共同下降时，一次查看、至多两次查看和及时查看对应不同数量级的额外资本；即使及时查看，锁定核验预算仍可能保留额外损失。正文只保留这两组比较及一个有限价格下的政策核查。研究对象是外部信号评估，并非运行自有策略来学习市场；可信样本偏差、固定市场及费用制度均需现实依据。结论是在规定制度中的理论比较，不构成投资建议或现实盈利验证。
\end{abstract}

\section{研究问题与决策者}
投资者准备使用一个筛选机会的信号，但不知道它是否准确。核验需要购买样本；决定采用后，还需为一次未来信号和实际经营支付费用。本文问：若核验与经营共用一笔初始资金，能否根据核验结果调整采购，会改变多少资金需求？

本文研究的投资者可以是个人或机构。信号提供者和评估使用者是不同角色；核验服务可以与信号提供者不同。本文不假设对方公开策略代码，也不把查看历史收益等同于取得可靠的未来表现证据。实际应用首先须说明核验样本从哪里来、真实类型如何查明、样本与未来机会的差异如何限定。以下将这些作为合同给定条件，尚未通过实际市场建立。付费抽样与最终实施收益共同分析已有直接研究\cite{wald}，本稿不将按下游收益选择信息本身作为一般方法首创。

流程被规定为先核验、通过后经营。若项目允许边运行自己的策略边学习，试验可能同时取得收入，信息也由投入产生，那是另一问题，不能直接套用本文资本结论。设置先核验的流程，需要内部风险要求或其他制度依据；本文比较该流程内部的决策权限，不证明它是所有投资任务的最佳流程。

\section{机会、核验服务与现金账本}
\subsection{经营收入怎样产生}
未来机会有快、慢两类，成功均支付1单位。捕获速度分别为2和0.01，快类份额 $u\in[0.5,0.6]$。若持续尝试到时间 $t$，速度为 $z$ 的机会成功概率为 $1-e^{-zt}$；成功即停止支出，否则到期退出。每单位时间成本 $c=0.2$，故扣除尝试成本的期望利润为
\begin{equation}
 F_z(t)=\left(1-\frac c z\right)(1-e^{-zt}).
\end{equation}
这里用到了平均尝试时长 $\int_0^t e^{-za}\,da$，不能把每次成本都算成 $ct$。$ct$ 是必须能支付的最长尝试预算。信号费 $f=0.001$ 另扣，尚未包括核验费用。以上金额均为模型单位。

信号把快类误判为慢类的概率为 $e_H$，反向错误为 $e_L$。合同以
\begin{equation}
 \theta=0.89e_H+0.11e_L
\end{equation}
评价质量。这是固定评价权重，不是未来类型比例。规定 $\theta\le0.01$ 的市场为好类，$\theta>0.05$ 的市场为不合格类。检验须在每个不合格市场中将通过概率控制在 $\alpha$ 以内，在每个好市场中将拒绝概率控制在 $\beta$ 以内，当前 $\alpha\approx8.06\times10^{-5}$、$\beta\approx3.78\times10^{-6}$。中间区域没有同样的通过率和利润承诺。固定质量权重与真实经济损失的关系尚需应用解释。

\subsection{样本可以买到什么，不能免费知道什么}
$p>0$ 是核验服务报价中的基本价格参数：服务商的各项价格都写成“固定倍数乘 $p$”。它由比较的报价情景给定，不是投资者要估计的市场参数，不是错误概率，也不是样本数量或样本精度。一次比较先固定 $p$，投资者再决定买哪种服务、买多少、何时停止。所有金额以一次捕获收入1为单位，因而以下对数中的 $p$ 是按该单位表示的价格数值。

令 $a_H,a_L\in(0,\infty)$ 分别表示保证快类、慢类样本的价格系数。二者是事先已知的报价条件，在 $p\downarrow0$ 时保持不变；可以取任意固定正数，不预设二者大小关系。固定精度 $\eta=\tau$ 时，两项报价分别为 $a_Hp$、$a_Lp$。例如，$a_H=2,a_L=20$、$p=2\times10^{-8}$ 时，二者分别收费 $4\times10^{-8}$ 和 $4\times10^{-7}$；该组数值只用于后文有限报价实例。将 $p$ 减半，所有服务的报价和退款都减半，但信号质量、核验样本偏差上限、未来经营成本、利润目标和可靠性要求不变。因此 $p\downarrow0$ 研究的是同一核验服务体系越来越便宜时的资本需求，而非好信号越来越少见或检验越来越准确。

自然抽样先预付 $10p$，交付类型为快类时退 $9p$，慢类不退款；慢类交付概率在 $[0.1,0.15]$ 内且事先固定。保证交付快、慢类型的价格分别为
\begin{equation}
 a_Hp(\tau/\eta)^r,\qquad a_Lp(\tau/\eta)^r,
 \quad0<\eta\le\tau=0.009,\quad r>0.
\end{equation}
保证服务的报价在样本交付前全额支付，没有退款；自然抽样仍按上述预付与退款制度结算。改变价格系数不改变服务的观测信息或质量承诺。对应样本的类型条件错误率与未来 $e_H,e_L$ 的差异至多为 $\eta$。这个已知差异范围是强假设，不能仅靠增加同源样本数量证明。零偏差没有有限报价。

市场参数和各服务的错误率在采购前固定。给定市场和所选服务后，样本具有新鲜独立随机性；未来经营机会与全部核验历史独立。不允许借款，闲置现金不计息，核验无收入回流。主要比较没有额外读取费或等待损失。修改这些条件，需要重做资金与统计推导。

\subsection{本金、费用与利润分别计算}
初始资金为 $B$，实际净核验费为 $C$，通过指示为 $I\in\{0,1\}$。通过后以余额 $B-C$ 买信号并选择尝试时长；不通过则退出。若经营收入为 $R$、经营支出为 $D$，退出时将二者记为零，则期末资金与完整利润分别为
\begin{align}
 W&=B-C-If-D+R,\\
 \Pi(B)&=\mathbb E(W-B)=\mathbb E[R-D-If-C].
\end{align}
未用本金返还，不作收入。要求每个好市场中 $\Pi(B)\ge g=0.1$，而且每笔预付款当时付得起。平均账单低不能替代付款保证；退款前余额也须足够。

\section{为什么要建立一个参数已知的对照}
\subsection{已知统计规律不等于知道交易结果}
为了分离信息不足带来的资金需求，先考虑一个理想对照：投资者免费、准确知道总体市场参数，不必核验，但仍不知道下一机会的真实类型和捕获结果，仍支付未来信号费与经营成本。

选定好市场 $G$：$u=0.5,e_H=1/89,e_L=0$。收到快标记的概率为 $44/89$，此时确为快机会。慢标记后的初始平均速度为 $2.89/90<c$，未成功后的平均速度继续下降，退出最优。快标记后的期望利润随尝试时长增加，因此用尽可支付的时长上限。资金为 $x\ge f$ 时，这个经营分支的期望净利润是
\begin{equation}\label{eq:v}
 v(x)=\frac{198}{445}\left[1-e^{-2(x-f)/c}\right]-f.
\end{equation}
它是先对每个可行经营方案求期望、再取最好方案的结果，不是对每次随机结果事后挑最幸运的选择，也不是每次保证赚到的金额。信号费与经营成本已扣，核验费未扣。

余额不足或经营无利可图时允许退出。完整经营价值为
\begin{equation}
 w(x)=\begin{cases}0,&0\le x<f,\\\max\{0,v(x)\},&x\ge f.\end{cases}
\end{equation}
后文的资本邻域中 $w=v>0$，才可直接用式\eqref{eq:v}。不能把该分支公式无条件外推到零余额。

\subsection{参照资本不是核验费}
记 $B_0$ 为上述参数已知对照达到目标 $g$ 的最低启动资金。由 $v(B_0)=g$ 得
\begin{equation}
 B_0=f-\frac c2\log\left(1-\frac{g+f}{198/445}\right)
 \approx0.02674696968.
\end{equation}
例如，此时最长时长约0.128735，快标记后的成功概率约0.226995，扣尝试成本后条件期望利润约0.204295；乘快标记概率得到0.101，再扣信号费得到0.1。仍有不成功而亏损的路径。$B_0$ 是固定市场 $G$ 中的经营资金基准，不是核验预算，更不是所有市场的同一最低资本。

\subsection{怎样定义最低资本及其增量}
先固定价格系数 $a_H,a_L$、价格 $p$ 和一种查看安排 $j$。下文的 $\mathcal A_j(B,p)$、$B_j^*(p)$ 均依赖这两个固定系数，记号中予以省略。一个完整方案不仅规定买多少样本，还规定选择哪种服务、如何依据已看到的结果继续或停止、何时通过，以及通过后如何用剩余资金经营。该方案不能预先知道真实市场参数。

记 $\mathcal A_j(B,p)$ 为满足下列全部要求的方案集合：采用安排 $j$ 允许的查看方式；在全部允许市场中按规定逐笔支付费用且不借款，退款前也能支付预付款；在每个不合格市场中通过概率不超过 $\alpha$；在每个好市场中拒绝概率不超过 $\beta$，且扣除核验费后的完整期望净利润至少为 $g=0.1$。市场范围包括前文允许的核验服务偏差，不能只对选定好市场 $G$ 达标。定义
\begin{equation}\label{eq:capital-definition}
 B_j^*(p)=\inf\{B\ge0:\mathcal A_j(B,p)\ne\varnothing\},
 \qquad \Delta_j^*(p)=B_j^*(p)-B_0.
\end{equation}
空集的下确界约定为 $+\infty$。使用下确界，是因为资金门槛可以被一系列方案逐渐逼近，而未必存在恰好用该金额的最优方案。若门槛可达到，才可以将它理解为某个方案实际使用的最小本金。

$\Delta_j^*(p)$ 就是本文的\textbf{最低资本增量}：在给定报价和查看权限下，达到全部要求所需的资金门槛，减去统一的经营基准 $B_0$。这里对全部允许方案取下确界，不是只优化某个固定检验的样本数量。基准 $B_0$ 来自固定市场 $G$，并非直接定义成全部好市场的已知参数门槛；本文另以全好类经营比较证明该基准可支持资本上界，不能仅凭命名就把两者等同。

例如，仅为说明定义，若某安排的门槛为0.026826，减去 $B_0\approx0.02674697$，其最低资本增量约为0.00007903。这个差额不是核验账单；它还受经营资金减少、错误退出和逐笔付款要求影响。正文后面的有限方案仅给最低门槛的上界，不应把该方案的资金直接称作已经求出的最小值。

对任意给定资金 $B$，实际资本与基准的差记为 $\delta_B=B-B_0$。只有对全部允许方案取下确界后，才得到 $\Delta_j^*(p)$。它不等于 $C$，也不意味着每条路径 $C\le\delta_B$。例如，只作账目演示，若 $B=B_0+0.0001$、检查必过且固定花费0.0001，经营仅剩 $B_0$，完整期望利润至多 $0.1-0.0001=0.0999$。把额外资本全部用于检查，不足以保留原利润目标。

\section{反馈怎样通过利润要求影响资本}
\subsection{同一个市场中的必要账目}
在固定 $G$ 中记实际拒绝概率为 $\varepsilon$。由于未来机会与核验历史条件独立，通过后经营价值不超过 $w(B-C)\le v(B)$，故
\begin{equation}\label{eq:easy}
 \Pi_G(B)\le(1-\varepsilon)v(B)-\mathbb E_G C.
\end{equation}
这不需要 $C$ 与通过事件独立。右侧仍乐观地允许通过后使用完整 $B$ 的经营能力，因此只是上界。要达到 $g=v(B_0)$，必须有
\begin{equation}\label{eq:margin}
 \varepsilon v(B)+\mathbb E_G C\le v(B)-v(B_0).
\end{equation}
左侧是收益损失和检查费用，右侧是新增资本产生的期望经营利润余量，两侧单位均为金额。由于 $v$ 在 $B_0$ 附近光滑且导数有限，式\eqref{eq:margin}迫使平均费用与实际拒绝概率至多为 $\delta_B$ 的固定常数倍。固定的合同上限 $\beta$ 本身并不要求拒绝率趋零；这里是利润目标在起作用。

更强的必要界来自对所有可能余额同时有效的经营切线上界。在该邻域内，$v(B)-Bv'(B)>0$，故连同退出选项也有 $w(x)\le v(B)+v'(B)(x-B)$，$0\le x\le B$。于是
\begin{equation}
 g\le v(B)-[1+v'(B)]\mathbb E_G C.
\end{equation}
这个加强项同时计入核验支出本身以及减少经营余额的影响。单凭 $v$ 的凹性而不检查退出分支，不能得到全余额比较。

\subsection{三种查看安排的资本比较}
一次查看指全部采购后才读取错误结果，此后不能追加采购；至多两次查看允许首批随机和根据第一次结果选择第二批；及时查看允许根据每次合法读取的结果决定后续动作。未读取或弃用的已付样本仍计费。

\begin{theorem}[同一任务下的查看时机与资本]
任取并固定 $a_H,a_L\in(0,\infty)$，保持上述服务信息、风险要求及利润目标 $g=0.1$ 不变，采用式\eqref{eq:capital-definition}的最低资本定义。当 $p\downarrow0$ 时，
\begin{align}
 B_{\rm one}^*(p)-B_0&=\Theta(p\log(1/p)),\\
 B_{\rm two}^*(p)-B_0&=\Theta(p\log\log(1/p)),\\
 B_{\rm online}^*(p)-B_0&=\Theta(p).
\end{align}
上下界的正常数及开始成立的价格阈值可依赖 $a_H,a_L$ 和其他固定合同参数，不声称对所有正价格系数同时一致。三种安排使用同一报价，不能通过改变检验要求或利润目标作比较。$\Theta$ 表示充分小价格下有正的固定倍数上下界，不是有限价格的准确报价。
\end{theorem}

\noindent\textbf{证明的连接。}
令要求的好侧拒绝率为 $\varepsilon\downarrow0$，固定坏侧上限。附录在一般价格系数下分别证明三类的平均核验费阶为 $p\log(1/\varepsilon)$、$p\log\log(1/\varepsilon)$ 与 $p$：必要方向在同一 $G$ 与一个固定不合格市场间作日志似然比较；充分方向用固定配额、两批筛选及可随时停止的证据规则，覆盖全部允许市场而非仅这一对。

式\eqref{eq:margin}给 $\varepsilon\le K\delta_B$、$\mathbb E_GC\le K\delta_B$，接统计必要界即得资本下界。充分方向取 $\varepsilon=p$；有限采购上限使最大账单及临时预付为 $O(p\log(1/p))\to0$，共享经营本金最终可逐笔支付。对实际账单调整经营时长，用经营函数的局部导数界控制期望利润损失，增加相应阶的资本即达原目标。完整统计引理、全好类经营比较及付款核对另附，不能仅凭一个期望费用界代替本证明。

\noindent\textbf{经济含义与边界。}
反馈的作用是把更多采购留给需要继续判断的路径，同时在明确的路径提前停止。最坏账单仍可带对数项；资本结论不等于每条路径都以平均账单付款。由于该最坏账单趋零而 $B_0>0$，主要渐近比较并非始终处于濒临破产的状态，而是费用、拒绝收益损失与经营资金共同作用。

\section{剩余核验资金能否重新使用}
在看到任何错误结果前，允许随机分出核验预算 $K$ 和经营预算 $X$，满足 $K+X=B$。如果未用 $K$ 不能转去经营，即使及时读取，对每个固定的 $a_H,a_L>0$ 仍有
\begin{equation}
 B^*_{\rm locked}(p)-B_0=\Theta(p\log(1/p)).
\end{equation}
若允许转出未用核验金，但每单位损失比例 $\phi\in[0,1]$，则对每个固定的 $a_H,a_L>0$，有对全部 $\phi$ 一致的正常数界：
\begin{equation}
 B^*(p,\phi)-B_0=\Theta\!\left(p[1+\phi\log(1/p)]\right).
\end{equation}
这项比较的必要方向用同一经营切线限制事前预留预算，再限制其能够支持的信息；充分方向用有限在线检查，完成后转入剩余预算。随机分账并不能绕过必要界。上下界常数仍可依赖这两个价格系数；详细不等式见附录。

其含义不是“账户越灵活越好”这一弱排序，而是剩余预算不能用于业务时，预留采购能力本身占用经营资源。该限制需要真实制度解释；若现实中原本可自由调拨，就应以灵活账户为主模型，不应人为锁定后再把解锁收益当成普遍现象。

\section{有限报价核查与适用范围}
正文只保留一项有限报价核查，不罗列研究轮次。仅在实例 $a_H=2,a_L=20$、$p=2\times10^{-8}$ 时，已有整数停止规则与事前随机组合给出充分资本小于0.026826086，最多13,920条标签，最坏净账单0.0007424。全部临时预付、退出费用、两侧检验要求和好侧利润均纳入验证。这些有限金额及标签数只适用于该实例，不直接用于其他价格系数。它是可行方案的保证，不是精确最优值；在某些好市场的平均费用并未同时改善。

本次在隔离副本实际重跑原有有理区间核验程序，重新检查七种报价中的资本、尾界与采购费用。没有重新运行全部候选搜索，也没有新增真实数据。数值核验支持具体有限不等式，渐近定理仍依靠解析证明。同一AI的再次核查不称为独立专家评审。

本次推广的是保证样本的固定正价格系数，并未同时推广利润目标、质量门槛或经营参数。若系数随 $p$ 趋零或发散，或降价同时改变样本质量，需要重新分析；不能从本定理推出同样结论。本文没有解决连续精度服务下的全局最优控制。有限菜单状态表示、精度离散近似及进一步证据增强移入技术附件，避免遮住资本比较主线。接下来的实质工作是建立外部核验服务及质量门槛的应用依据，并检验读取延迟和样本迁移；如果要回到自有策略，则必须重新建立试运行收入、观察与资金回流模型。

\appendix
\section{核验费用的三个统计量级}
本附录将正文需要的统计比较集中说明，避免将样本费用直接称为资本。固定 $\alpha<1/2$，只令好侧实际拒绝要求 $\varepsilon\downarrow0$。全部常数可依赖既定质量间隔与固定价格系数 $a_H,a_L$。为核对每个费用界，记
\begin{equation}
 b_H=\min\{1,a_H\}>0,\qquad
 b_*=\min\{1,a_H,a_L\}>0.
\end{equation}
自然抽样交付快类的净价为 $p$，保证快类服务净价至少为 $a_Hp$，所以每个已付快标签至少净花 $b_Hp$。任何类型的已付标签至少净花 $b_*p$。这些是净费用下界；预付款是否付得起仍单独检查。

\subsection{一对固定市场给必要费用}
取正文 $G$，再取固定不合格市场 $Q$：$e_H=9/89,e_L=0$。在两市场中所有服务都准确对应未来错误率，故满足全部精度合同；类型供应与付款元数据分布相同。每个已付快类标签至少花 $b_Hp$。定义 $D=\operatorname{kl}(9/89,1/89)>0$。

一次查看时，全部采购在错误可见前完成，两市场的购买数分布相同。完整记录反向相对熵至多 $D\mathbb E_GC/(b_Hp)$。不合格市场拒绝至少 $1-\alpha$，好市场拒绝至多 $\varepsilon$，合并到拒绝事件给
\begin{equation}
 \frac{D\mathbb E_GC}{b_Hp}\ge\operatorname{kl}(1-\alpha,\varepsilon)
 \ge(1-\alpha)\log(1/\varepsilon)-\log2.
\end{equation}
因此费用至少为 $p\log(1/\varepsilon)$ 阶。精度选择、额外付费后弃用和私人随机采购不减少此必要界。

至多两次查看时，令 $M=\mathbb E_GC/(b_Hp)$，$K$ 为第一次错误可见前全部已付快标签数。两市场的 $K$ 同分布，且 $\mathbb E K\le M$，所以共同事件 $A=\{K\le2M\}$ 的概率至少一半。条件于 $A$，第一次记录的 $Q/G$ 密度比至多 $9^{2M}$。第一次记录相同后的第二批动作规则相同，但其无条件期望购买数可以不同；用上述密度比转换，得
\begin{equation}
 \operatorname{KL}(Q_{\rm rec\mid A},G_{\rm rec\mid A})
 \le2DM9^{2M}.
\end{equation}
条件错误概率至多为原来的两倍，故充分小 $\varepsilon$ 时
\begin{equation}
 2DM9^{2M}\ge\operatorname{kl}(1-2\alpha,2\varepsilon).
\end{equation}
右侧至少为 $\log(1/\varepsilon)$ 的固定正常数倍；两侧取对数，得到 $M=\Omega(\log\log(1/\varepsilon))$。这里没有把反馈后的两个市场平均样本数设成相同。

及时查看至少以概率 $1-\alpha-\varepsilon$ 实际付款：否则在没有任何付费结果的记录上，两市场采取相同决定，无法形成所需的通过概率差。因此 $\mathbb E_GC\ge b_*p(1-\alpha-\varepsilon)$。

\subsection{一组简单可行规则给对应上界}
只使用固定精度 $\eta=\tau$ 的保证服务，每批买一个快标签和一个慢标签，价格 $(a_H+a_L)p$。令二元错误为 $X_H,X_L$，用 $Y=0.89X_H+0.11X_L\in[0,1]$。任何好市场的批均值至多 $a=0.019$，任何不合格市场至少 $d=0.041$。取中点 $m=0.03$、$h=m-a=d-m=0.011$。各批独立，以下Hoeffding界覆盖全部允许服务偏差。

一次查看取固定 $N$，批均值不超过 $m$ 就通过。好拒绝与坏通过分别不超过 $e^{-2Nh^2}$。取 $N$ 同时满足两侧要求，数量为 $O(\log(1/\varepsilon))$。

两次查看先设末批量 $N_2=O(\log(1/\varepsilon))$，令其好拒绝至多 $\varepsilon$、坏通过至多 $\alpha/2$。首批 $N_1$ 同时使坏通过至多 $\alpha/2$、好侧未通过概率至多 $1/N_2$，故 $N_1=O(1+\log N_2)$。首批通过即结束，否则购买新鲜末批并按其结果决定。坏通过总概率至多 $\alpha$，好拒绝至多 $\varepsilon$；好侧平均批数至多 $N_1+1$，最大批数 $N_1+N_2$。

及时查看令累积对数证据为 $S_n=k\sum_{i=1}^n(m-Y_i)$，取 $k=4h$。坏市场中
\begin{equation}
 \mathbb E e^{k(m-Y)}\le e^{k(m-d)+k^2/8}\le1.
\end{equation}
所以 $S_n$ 首达 $\log(1/\alpha)$ 的概率至多 $\alpha$。好市场的平均增量至少 $kh>0$。设置有限上限
\begin{equation}
 N\ge\max\left\{\frac{2\log(1/\alpha)}{kh},
                    \frac{2\log(1/\varepsilon)}{h^2}\right\}
\end{equation}
并在上限未通过时拒绝，Hoeffding界给拒绝概率至多 $\varepsilon$。令 $\sigma$ 为通过与上限的较早时刻，有限停止的漂移恒等式与正越界至多 $km$ 给
\begin{equation}
 kh\mathbb E_G\sigma\le\mathbb E_G S_\sigma
 \le\log(1/\alpha)+km.
\end{equation}
故平均批数一致有界，最大批数仍为 $O(\log(1/\varepsilon))$。这组简单规则只证明量级，不用于正文有限报价纪录。与原数值系数相比，统计量和样本数的概率界均未改变；平均费用只乘以固定系数 $a_H+a_L$。保证样本逐笔预付，在至多 $N_{\max}$ 批的程序中，累计支出及付款峰值不超过 $(a_H+a_L)pN_{\max}$。上述三种程序均有 $N_{\max}=O(\log(1/\varepsilon))$，所以取 $\varepsilon=p$ 时该峰值趋零。对固定的两个价格系数，这给出了逐笔可支付的上界，而非仅给平均费用。

\section{从好市场的经营收益到资本上界}
正文采用只经营快标记的规则。快类份额为 $u$ 时，余额决定时长 $t$，经营期望为
\begin{equation}
 R(e_H,e_L,u;t)=u(1-e_H)F_2(t)+(1-u)e_LF_{0.01}(t)-f.
\end{equation}
由于 $F_2\ge0,F_{0.01}\le0$，最坏 $u=0.5$。对好三角形 $0.89e_H+0.11e_L\le0.01$，该式仿射，最小值在三个顶点中取得。$0<t\le0.15$ 时，
\begin{equation}
 11F_2(t)\ge9.9(1.72775)t>16.91t\ge-89F_{0.01}(t).
\end{equation}
第一步用 $1-e^{-2t}\ge2t-2t^2+4t^3/3-2t^4/3$，最后一步用 $1-e^{-t/100}\le t/100$。因此最差顶点是 $(e_H,e_L)=(1/89,0)$，经营期望至少为正文 $v(B-C)$。该证明只覆盖规定时长范围，不外推至所有余额。

取 $\varepsilon=p$。若最大账单和预付峰值为 $K_p=O(p\log(1/p))$，则充分小 $p$ 时，全部余额 $B-C$ 仍在 $B_0$ 附近的正值区间。设此区间的导数上下界为 $0<m_0\le v'\le L_0$，则无须通过与费用独立，也有
\begin{equation}
 \Pi(B)\ge(1-p)v(B)-(1+L_0)\mathbb EC.
\end{equation}
对各统计费用函数 $f_j(p)$，取 $B=B_0+Kpf_j(p)$，导数下界给新增经营值至少 $m_0Kpf_j(p)$。选充分大固定 $K$，可覆盖拒绝损失和全部平均费用。有限最大账单负责逐笔支付，平均账单负责期望利润，两者在证明中承担不同作用。

\section{分账与转账的必要方向}
设事前核验预算为 $K$、实际账单为 $C\le K$。即使把所有未用预算转出，经营余额至多 $B-\phi K-(1-\phi)C$。为给政策更有利的上界，可暂时忽略转账损耗本身的利润扣除；经营切线仍给
\begin{equation}
 v'(B)\phi\mathbb E_GK+
 [1+v'(B)(1-\phi)]\mathbb E_GC\le v(B)-g.
\end{equation}
每条路径最多采购 $K/(b_Hp)$ 个快标签，$K$ 在两困难市场中的分布相同，反向日志信息因此至多 $D\mathbb EK/(b_Hp)$。事件错误界迫使 $\mathbb EK$ 至少为 $p\log(1/\delta_B)$ 阶；另有 $\mathbb EC\ge c_1p$，其中 $c_1>0$ 为固定常数。与切线联合，即得统一必要资本阶。充分方向用前述在线规则、有限预留预算和实际剩余资金转入经营，转账损失及余额减少均为 $O(\phi p\log(1/p)+p)$。这里的正常数可依赖固定价格系数，关于 $\phi$ 的一致性保持不变。本论证解释量级；有限政策的每笔退款、尾界与实际金融验收仍见随附原核验记录。

\begin{thebibliography}{9}
\bibitem{wald} Adusumilli K. How to Sample and When to Stop Sampling: The Generalized Wald Problem and Minimax Policies. 2025，2025年3月25日作者稿。核验摘要和引言，以实验成本与后续实施收益的联系为背景；本文资本阶不由该文直接推出。\href{https://www.restud.com/wp-content/uploads/2025/03/Sequential_stopping_final.pdf}{原文}。
\end{thebibliography}
\end{document}
